% !tp-example Basic usage
\documentclass{article}
\begin{document}
\wcTpTheorem[type=theorem,heading=Nonnegativity,number=1,id=square]{
  The square of every real number is nonnegative.

}
\end{document}

% !tp-example Javascript API
\documentclass{article}
\begin{document}
\wcScript[type=module]|
import "/components/theorem/theorem.js";

const example = document.createDocumentFragment();

const element1 = document.createElement("tp-theorem");
element1.setAttribute("type", "theorem");
element1.setAttribute("heading", "Nonnegativity");
element1.setAttribute("number", "1");
element1.setAttribute("id", "square");
element1.append(document.createTextNode("\n  "));

const element2 = document.createElement("p");
element2.append(document.createTextNode("The square of every real number is nonnegative."));
element1.append(element2);
element1.append(document.createTextNode("\n"));
example.append(element1);

document.body.append(example);
|
\end{document}

% !tp-example Statement types
\documentclass{article}
\begin{document}
\wcTpTheorem[type=theorem]{
  A theorem with its usual label.

}

\wcTpTheorem[type=lemma]{
  A lemma with its usual label.

}

\wcTpTheorem[type=proposition]{
  A proposition with its usual label.

}

\wcTpTheorem[type=corollary]{
  A corollary with its usual label.

}

\wcTpTheorem[type=definition]{
  A definition with its usual label.

}

\wcTpTheorem[type=example]{
  A example with its usual label.

}

\wcTpTheorem[type=remark]{
  A remark with its usual label.

}

\wcTpTheorem[type=proof,qed=□]{
  A proof with its usual label.

}
\end{document}

% !tp-example Shared section counters
\documentclass{article}
\begin{document}
\section*{First section}

\wcTpTheorem[numbering=section,counter=results]{
  The first statement.

}

\wcTpTheorem[type=lemma,numbering=section,counter=results]{
  A supporting lemma.

}

\section*{Second section}

\wcTpTheorem[numbering=section,counter=results]{
  The counter restarts.

}
\end{document}

% !tp-example Background variants
\documentclass{article}
\begin{document}
\wcTpTheorem[type=remark,numbering=none,variant=none,heading=none]{
  The default appearance has no border or background.

}

\wcTpTheorem[type=remark,numbering=none,variant=neutral,heading=neutral]{
  The background follows TP theme tokens.

}

\wcTpTheorem[type=remark,numbering=none,variant=brand,heading=brand]{
  The background follows TP theme tokens.

}

\wcTpTheorem[type=remark,numbering=none,variant=info,heading=info]{
  The background follows TP theme tokens.

}

\wcTpTheorem[type=remark,numbering=none,variant=success,heading=success]{
  The background follows TP theme tokens.

}

\wcTpTheorem[type=remark,numbering=none,variant=warning,heading=warning]{
  The background follows TP theme tokens.

}

\wcTpTheorem[type=remark,numbering=none,variant=danger,heading=danger]{
  The background follows TP theme tokens.

}
\end{document}

% !tp-example Ending symbols
\documentclass{article}
\begin{document}
\wcTpTheorem[type=theorem,qed=◆]{
  A theorem with a diamond at its end.

}

\wcTpTheorem[type=definition,qed=∎]{
  A definition with a filled square.

}

\wcTpTheorem[type=proof,qed=□]{
  A proof with an explicit closing square.

}

\wcTpTheorem[type=proof,qed]{
  The explicit empty value hides the closing square.

}
\end{document}

% !tp-example Rich heading and citation
\documentclass{article}
\begin{document}
\wcDiv[data-tp-reference-scope]{
  \wcTpTheorem[type=proposition,number=2.1]{
    \wcSpan*[slot=heading]{Companion input, \wcTpRef*[href=@guide]{}}
    The conclusion follows from the cited result.

  }
  \rawhtml|<tp-biblio ref="guide" label="[1]"><p>Alex Example. <em>A guide to reconstruction</em>. 2026.</p></tp-biblio>|
  \wcTpListof[selector=tp-biblio]{}
}
\end{document}
